\chapter{Self Adjoint Operators}
\section{Inner Product Space}
Let $K = \RR$ or $K =\CC$.
\begin{enumerate}
    \item When $K = \RR$, for $a\in \RR$, $|a| =\max \{a, -a\} \in \RR_{\ge 0}$.
    \item When $K=\CC$, for $z = x+ \ii y \in \CC$, $x,y\in \RR$, $|z|= \sqrt{x^2 +y^2} \in \RR_{\ge 0}$.
\end{enumerate}
\begin{definitionenv}
    We fix a vector space over $K$, $\langle \cdot,\cdot\rangle : V\times V \rightarrow K$.
    We say that $\langle\cdot,\cdot\rangle$ is a semi-inner product if the following conditions hold:
    \begin{enumerate}
        \item $\langle \cdot ,\cdot\rangle$ is linear with respect to the second variable, namely for $(y,z) \in V\times V$, $(a,b) \in K\times K$, $x\in V$,
        $$ \langle x, ay+bz\rangle = a\langle x,y\rangle + b\langle x,z \rangle.$$
        \item $\langle \cdot ,\cdot\rangle$ is conjugate symmetric, namely for $(x,y)\in V^2$, 
        $$ \langle y ,x\rangle = \overline{\langle x,y\rangle}.$$
        \item $\langle \cdot ,\cdot\rangle$ is semi-positive, namely for $x\in V$,
        $$ \langle x,x\rangle \ge 0.$$
    \end{enumerate}
    The pair $\left(V,\langle \cdot ,\cdot\rangle\right)$ is called \textbf{semi-inner product space}.
    When $K=\RR$ a semi-inner product on $V$ is just a positive semi-definite bilinear form. If in addition, $\langle x,x\rangle>0$ when $x\in V\setminus \{0\}$, we say that $\langle \cdot ,\cdot\rangle$ is an \textbf{inner product} and $\left(V,\langle \cdot ,\cdot\rangle\right)$ is an \textbf{inner product space}.
\end{definitionenv}
\begin{remark}
    Let $\left(V,\langle \cdot ,\cdot\rangle\right)$ be a semi-inner product space, $(x,y,z)\in V^3$, $(a,b)\in K^2$.
    $$ \langle ax+by,z\rangle = \overline{\langle z, ax + by\rangle} = \overline{a\langle z,x\rangle + b\langle z,y\rangle} = \overline{a} \langle x,z \rangle +\overline{b} \langle y,z \rangle.$$ 
\end{remark}
\begin{exampleenv}
    Let $n\in \NN_{> 0}$, $\langle \cdot ,\cdot\rangle : K^n \times K^n \rightarrow K$,
    $$ \langle x,y\rangle = \sum_{i=1}^{n} x_i \overline{y_i}$$
    is an inner product. We call it the \textbf{canonical inner product}.
\end{exampleenv}
\begin{definitionenv}
    Let $\left(V,\langle \cdot ,\cdot\rangle\right)$ be a semi-inner product space. If $x,y\in V$ are such that $\langle x,y\rangle =0$, then we say that $x$ and $y$ are \textbf{orthogonal} and we write $x\perp y$.

    \quad Note that if $x$ and $y$ are orthogonal, then for any $(a,b)\in K^2$, 
    $$ \langle ax , by\rangle = \overline{a}b \langle x, y\rangle =0.$$
    Moreover, 
    $$ \langle x+y, x+y \rangle = \langle x,x \rangle + \langle y,y\rangle.$$
    For any $A\subseteq V$, we denote by $A^\perp$ the set
    $$ A^\perp \coloneq \{x\in V\mid \forall y\in A, \langle x,y\rangle =0 \}.$$
    $V^\perp$ is called the \textbf{kernel} of the semi-inner product.
\end{definitionenv}
\begin{propositionenv}\label{3.1.5}
    Let $\left(V, \langle \cdot, \cdot\rangle\right)$ be a semi-inner product space.
    \begin{enumerate}
        \item For any $A\subseteq V$, $A^\perp$ is a vector subspace of $V$.
        \item For any $A\subseteq V$, the vector subspace of $V$ generated by $A$ is contained in $\left(A^\perp\right)^\perp$.
        \item For any $A,B\subseteq V$, if $A\subseteq B$, then $B^\perp \subseteq A^\perp$.
        \item For any $A\subseteq V$, let $W$ be a vector subspace generated by $A$, then $$A^\perp = W^\perp.$$
    \end{enumerate}
\end{propositionenv}
\begin{proofenv}
    \quad 
    \begin{enumerate}
        \item Let $(x,y)\in A^\perp\times A^\perp$, $(a,b) \in K^2$, for any $z\in A$, $$ \langle ax + by, z \rangle = \overline{a} \langle x,z\rangle + \overline{b} \langle y,z\rangle =0.$$
        \item By 1. we only need to show $A\subseteq \left(A^\perp\right)^\perp$. For any $x\in A$, $y\in A^\perp$, we have $$ \langle x,y\rangle = \overline{\langle y , x\rangle} = 0,\ x\in \left(A^\perp\right)^\perp.$$
        \item Let $y\in B^\perp$. For any $x\in A$, we have $x\in B$, so $\langle y, x\rangle =0$, so $y\in A^\perp$.
        \item By 3. $W^\perp \subseteq A^\perp$, by 2. $W\subseteq \left(A^\perp\right)^\perp$. Therefore, by 2. and 3. $$ A^\perp \subseteq \left(\left(A^\perp\right)^\perp\right)^\perp \subseteq W^\perp.$$
    \end{enumerate}
\end{proofenv}
\begin{theoremenv}[Cauchy-Schwarz]
    Let $\left(V,\langle \cdot, \cdot \rangle\right)$ be a semi-inner product space.
    For any $(x,y)\in V^2$, one has 
    $$ \left| \langle x,y\rangle\right|^2 \le \langle x, x\rangle \cdot \langle y,y\rangle.$$
\end{theoremenv}
\begin{proofenv}
    Let $u\in K$ such that $|u|=1$, such that
    $$ u\langle x,y\rangle = \left|\langle x,y\rangle\right|.$$
    For any $t\in \RR$, one has 
    $$ 0\le \langle x + ty, x+ty\rangle = \langle x,x\rangle + 2t \left| \langle x,y\rangle \right| + t^2\langle y,y\rangle.$$
    So the discriminant of the quadratic function is non-positive.
\end{proofenv}
\begin{propositionenv}
    Let $\left(V, \langle \cdot, \cdot\rangle\right)$ be a semi-inner product space.
    \begin{enumerate}
        \item The mapping $$\begin{array}{rrcl}
            \|\cdot\|: & V & \rightarrow &\RR_{\ge 0}\\
            & x &\mapsto & \sqrt{\langle x,x \rangle}
        \end{array}$$
        is a semi-norm on $V$. We call it the \textbf{semi-norm associated with the semi-inner product} $\langle \cdot ,\cdot\rangle$.
        It is a norm when $\langle \cdot ,\cdot\rangle$ is an inner product.
        \item The kernel of $\langle \cdot ,\cdot\rangle$ is equal to $\{x\in V\mid \| x\| = 0\}.$
        \item Let $W$ be a vector subspace of $V$, such that $W\subseteq V^\perp$. Then the semi-inner product $\langle \cdot,\cdot\rangle$ induces a semi-inner product:
        $$\begin{array}{rrcl}
            \langle \cdot ,\cdot\rangle^\sim : & \left(V/W\right) \times \left(V/W\right) & \rightarrow & K\\
            & \langle [x], [y] \rangle & \mapsto & \langle x,y\rangle.
        \end{array}$$
        Moreover, if $W = V^\perp$, then $\langle \cdot, \cdot\rangle^\sim$ is an inner product on $V/W$.
    \end{enumerate}
\end{propositionenv}
\begin{proofenv}
    \quad 
    \begin{enumerate}
        \item Take any $(a,x) \in K\times V$, we have 
        $$\| ax\| = \langle ax, ax\rangle^{\frac{1}{2}} = |a| \langle x,x\rangle^{\frac{1}{2}} = |a|\|x\|.$$
        Take any $(x,y)\in V^2$, we have 
        $$ \| x+y\|^2 = \langle x+y,x+y\rangle = \langle x,x\rangle + 2\mathrm{Re}\langle x, y\rangle + \langle y, y\rangle.$$
        By Cauchy-Schwarz, $\mathrm{Re}\langle x,y\rangle \le |\langle x, y\rangle| \le \|x\|\|y\|$, hence
        $$ \| x+y \|^2 \le \| x\|^2 +\|y \|^2 + 2\|x\|\|y\| = \left(\|x\| +\|y\|\right)^2.$$
        \item $$ V^\perp \subseteq \{x \in V\mid \| x\| = 0\}.$$
        It remains to show the opposite inclusion. Let $x\in V$ such that $\| x\| = 0$. For any $y\in V$, by Cauchy-Schwarz,
        $$ |\langle x,y\rangle|^2 \le \| x\|^2 \cdot \| y\|^2 =0.$$
        So $\langle x,y\rangle =0$, so $x\in V^\perp$.
        \item Let $(x,y)\in V^2$, $(u,v)\in W^2$, we have
        $$ \langle x + u, y + v \rangle = \langle x,y\rangle + \langle x,v\rangle + \langle u,y\rangle + \langle u,v\rangle = \langle x,y\rangle.$$
        Hence, by passing to the quotient, $\langle \cdot, \cdot\rangle^\sim$ induces 
        $$\begin{array}{rrcl}
            \langle \cdot ,\cdot\rangle^\sim : & \left(V/W\right) \times \left(V/W\right) & \rightarrow & K\\
            & \left( [x], [y] \right) & \mapsto & \langle x,y\rangle.
        \end{array}$$
        We check $\langle \cdot ,\cdot\rangle^\sim$ is an inner product. Let $(x,y,z)\in V^3$, $(a,b)\in K^2$, we have
        $$ \begin{array}{rl}
            \langle [x] , a[y] + b[z]\rangle &= \langle x, a[y] + b[z]\rangle  \\
            &= a\langle x,y\rangle + b\langle x,z\rangle = a\langle [x], [y]\rangle^\sim + b\langle [x], [z]\rangle^\sim.
        \end{array}$$
        $$ \langle [y],[x] \rangle^\sim = \langle y,x\rangle = \overline{\langle x,y\rangle} = \overline{\langle [x], [y]\rangle^\sim}.$$
        $ \langle [x], [x]\rangle^\sim = \langle x,x\rangle \ge0.$ So $\langle \cdot, \cdot\rangle^\sim$ is a semi-inner product on $V/W$.
        If $W = V^\perp$, then $\langle [x], [x]\rangle^\sim = \| x\| = 0$ if and only if $[x]=0$ in $V/W$.
    \end{enumerate}
\end{proofenv}
\begin{propositionenv}\label{prop:3.1.8}
    Let $(V,\langle \cdot ,\cdot\rangle)$ be a inner product space. $\| \cdot \|: V\rightarrow \RR_{\ge 0}$ a norm such that $\| x\| = \sqrt{\langle x,x\rangle}$ for any $x\in V$.
    \begin{enumerate}
        \item For any $ x\in V $, the mapping
        $$\begin{array}{rrcl}
            l_x: & V & \rightarrow & K\\
            & y & \mapsto & \langle x,y\rangle
        \end{array}$$
        is bounded $K$-linear mapping whose norm is equal to $\|x\|$.
        \item For any $(x,y)\in V^2$, $l_x = l_y$ if and only if $x = y$.
    \end{enumerate}
\end{propositionenv}
\begin{proofenv}
    \quad
    \begin{enumerate}
        \item By Cauchy-Schwarz, $\left| \langle x,y\rangle \right| \le \left\| x \right\| \cdot \left\| y \right\|$, for any $y \in V$.
        So $l_x$ is bounded, and since $l_x (x) = \left\| x\right\|^2$ so the norm of $l_x$ is equal to $\| x\|$.
        \item One has $l_x - l_y = l_{x-y}$, so $\| l_x - l_y \| = \|x -y\|$. So $l_x - l_y =0$ if and only if $x - y =0$.
    \end{enumerate}
\end{proofenv}
\begin{propositionenv}\label{prop:3.1.9}
    Let $\left(V,\langle \cdot ,\cdot\rangle\right)$ be an inner product space, $\left\| \cdot \right\|$ a norm such that $\| x\| = \sqrt{\langle x,x\rangle}$ for any $x\in V$.
    \begin{enumerate}
        \item For any subset $A \subseteq V$, $A^\perp$ is a closed subset of $V$.
        \item Let $A \subseteq V$, $B = \bar{A}$, then $A^\perp = B^\perp$.
    \end{enumerate}
\end{propositionenv}
\begin{proofenv}
    \quad 
    \begin{enumerate}
        \item $$A^\perp = \bigcap_{x\in A} \ker\left(l_x\right).$$
        \item By proposition \ref{3.1.5}, $B^\perp \subseteq A^\perp$. To show the opposite inclusion, take $ y\in A^\perp$.
        We see that $A \subseteq \ker\left(l_y\right)$, since $l_y$ is continuous, $\ker\left(l_y\right)$ is closed. 
        Hence $B =\bar{A} \subseteq \ker\left(l_y\right)$.
    \end{enumerate}
\end{proofenv}
\begin{propositionenv}
    Let $\left(V,\langle \cdot ,\cdot\rangle\right)$ be an inner product space over $K$.
    The norm $\| \cdot \|$ is associated with some inner product $\langle \cdot ,\cdot\rangle$ if and only if for any $(x,y)\in V^2$,
    $$ \|x+y\|^2 + \|x-y\|^2 = 2\|x\|^2 + 2\|y\|^2. $$
\end{propositionenv}
\begin{proofenv}
    Suppose that there exist an inner-product such that $\| x\| = \sqrt{\langle x,x\rangle}$ for any $x\in V$. For any $(x,y)\in V^2$, we have
    $$ \| x+y\|^2 + \| x-y\|^2 = \langle x+y, x+y\rangle + \langle x-y, x-y\rangle = 2\| x\|^2 + 2\| y\|^2.$$
    Conversely, suppose that for any $(x,y)\in V^2$, we have
    $$ \| x+y\|^2 + \| x-y\|^2 = 2\| x\|^2 + 2\| y\|^2.$$
    We define $\langle \cdot ,\cdot\rangle : V\times V \rightarrow K$ as follows:
    $$ \langle x,y\rangle = \frac{1}{2}\left(\|x+y\|^2 - \|x-y\|^2\right).$$
\end{proofenv}

\section{Orthogonal Projection}
\begin{definitionenv}
    Let $\left(V, \langle \cdot, \cdot \rangle\right)$ be an inner product space. 
    If $V$ is complete with respect to the topology induced by the norm
    $$\begin{array}{rrcl}
        \left\| \cdot \right\|: & V & \rightarrow & \RR_{\ge 0}\\
        & x & \mapsto & \sqrt{\langle x,x\rangle},
    \end{array}$$
    then we say that $V$ is a \textbf{Hilbert space}.

    \quad In particular, if $V$ is finite-dimensional, then $V$ is a Hilbert space.
\end{definitionenv}

\begin{theoremenv}\label{3.2.2}
    Let $\left(V,\langle \cdot,\cdot\rangle\right)$ be a inner product space and $W$ be a vector subspace of $V$. We assume that $\left(W, \left.\langle\cdot,\cdot\rangle\right|_{W}\right)$ is a Hilbert space.
    
    \quad For any $x\in V$, there exist a \textit{unique} element $p_W(x) \in W$ such that $\left\| x - p_W(x)\right\|$ minimizes the function $ \left(y \in W\right) \mapsto \left\| x - y\right\|$.
    
    \quad $p_W(x)$ is the only element of $W$ such that $x - p_W(x) \in W^\perp$.
\end{theoremenv}
\begin{proofenv}
    We first prove the existence of $p_W(x)$. 
    Consider $\left( y \in W\right) \mapsto \left\| x- y\right\|$.
    Put $\dis c\coloneq \inf_{y \in W} \left\| x- y\right\|$.
    Take a sequence $\left(y_n\right)_{n\in \NN}$ in $W$ such that $$\lim_{n\rightarrow \infty} \left\| x - y_n\right\| = c.$$
    Note that for any $m,n\in \NN$,
    $$\begin{array}{rl}
        \left\| y_n - y_m\right\|^2 =& 2\left\| y_n - x\right\|^2 + 2\left\| y_m - x\right\|^2 - \left\| y_m + y_n -2x\right\|^2 \\
        = & 2\left\| y_n - x\right\|^2 + 2\left\| y_m - x\right\|^2 - 4 \left\| \frac{y_m + y_n}{2} - x \right\|^2 \\
        \le& 2 \left\| y_n - x\right\|^2 + 2 \left\| y_m - x\right\|^2 - 4 c^2.
    \end{array}.$$
    So, 
    $$ \lim_{ N \rightarrow +\infty} \sup_{n,m\ge N} \left\| y_n - y_m\right\|^2 = 0,$$
    and hence $\left(y_n\right)_{n\in \NN}$ is a Cauchy sequence in $W$.
    Put $\dis p_W(x) \coloneq \lim_{n\rightarrow \infty} y_n$, then
    $$ \left\| x - p_W(x)\right\| = \inf_{y\in W} \left\| x - y\right\|.$$
    We now prove that $x - p_W(x) \in W^\perp$.
    Let $y\in W$, $a\in K$ such that $|a|=1$, and $a \langle x -p_W(x) , y\rangle = \left| \langle x- p_W(x), y\rangle\right|.$
    For any $t\in \RR$,
    $$ \begin{array}{rl}
        0 \le& \left\| x - p_W(x) + t a y\right\|^2 \\
        =& \left\| x - p_W(x)\right\|^2 + 2t \left| \langle x - p_W(x), y\rangle\right| + t^2 \left\| y\right\|^2\\
        \ge& \left\| x - p_W(x)\right\|^2.
    \end{array}$$
    We deduce that $\forall t \in \RR$,
    $$2t\left| \langle x - p_W(x), y\rangle\right| + t^2 \left\| y\right\|^2 \ge 0,$$
    hence $\langle x - p_W(x) ,y\rangle =0$.

    \quad ``uniqueness'': Take any $z \in W$, 
    $$ \begin{aligned}
        \left\| x - z\right\| = \left\| (x - p_W(x)) + (p_W(x) - z)\right\|^2 =& \left\| x - p_W(x)\right\|^2 + \left\| p_W(x) - z\right\|^2 \\
        \ge& \left\| x - p_W(x)\right\|^2.
    \end{aligned}$$
    If $x - z \in W^\perp$,
    $$ \left\| x - z\right\|^2 = \left\| x - p_W(x)\right\|^2 + \left\| p_W(x) - z\right\|^2.$$
    So $p_W(x) = z$.
\end{proofenv}

\begin{definitionenv}
    Let $\left(V, \langle \cdot, \cdot\rangle\right)$ be an inner product space, $W$ a vector subspace of $V$ such that $\left(W, \left.\langle\cdot,\cdot\rangle\right|_W\right)$ is a Hilbert space. The mapping
    $$\begin{array}{rrcl}
        p_W: & V & \rightarrow & W\\
        & x & \mapsto & p_W(x)
    \end{array}$$
    is called the \textbf{orthogonal projection} from $V$ to $W$. $p_W(x)$ is called the \textbf{orthogonal projection} of $x$ to $W$.
\end{definitionenv}
\begin{definitionenv}
    Let $\left(V, \langle \cdot, \cdot\rangle\right)$ be an inner product space.
    
    \quad We say that a family $\left(x_i\right)_{i\in I}$ of non-zero elements in $V$ is \textbf{orthogonal} if for any $\left(i,j\right)\in I^2$, $i\neq j$, $\langle x_i, x_j\rangle = 0$. 
    
    \quad If additionally, if $ \left\| x_i \right\| = 1 $ for any $i\in I$, then we say that $\left(x_i\right)_{i\in I}$ is an \textbf{orthonormal} family.
    
    \quad If $\left(x_i\right)_{i\in I}$ is an orthogonal family such that  $\mathrm{Vect}_{K}\left(\left\{ x_i \mid  i\in I \right\}\right)$ is dense in $V$, then we say that $\left(x_i\right)_{i\in I}$ is an \textbf{orthogonal base} of $V$.
\end{definitionenv}
\begin{propositionenv}
    Let $\left(V, \langle \cdot, \cdot\rangle\right)$ be an inner product space.
    Any orthogonal family in $V$ is linearly independent in $V$.
\end{propositionenv}

\begin{proofenv}
    Without loss of generality, let $I = \left\{1,\cdots,n\right\}$, $\left(a_1,\cdots,a_n\right)\in K^n$ such that
    $$ a_1 x_1 + \cdots + a_n x_n =0.$$
    Then, for any $i\in I$, 
    $$ 0 = \langle a_1 x_1 + \cdots + a_n x_n, x_i\rangle = a_i \langle x_i, x_i\rangle.$$
    So $a_i =0$ for any $i\in I$, as $\langle x_i,x_i\rangle = \left\| x_i \right\|^2 \neq 0$.
\end{proofenv}
\begin{propositionenv}
    Let $\left(V, \langle \cdot, \cdot\rangle\right)$ be a Hilbert space, $\left(x_i\right)_{i\in I}$ orthogonal family in $V$. $\left(a_i\right)_{i\in I} \in K^I$. Suppose that 
    $$ \sum_{i\in I} \left|a_i\right|^2 \cdot \left\|x_i\right\|^2 \coloneq \sup_{\substack{J\subseteq I\\ J \text{finite}}} \sum_{i\in J} \left|a_i\right|^2 \cdot \left\|x_i\right\|^2 < +\infty.$$
    \begin{enumerate}
        \item Then the set $I' = \left\{i\in I \mid a_i \neq 0\right\}$ is countable.
        \item In the case when $I'$ is infinite, for any bijection $ \varphi : \NN \rightarrow I'$, the series
        $$ \sum_{n\in \NN} a_{\varphi(n)} x_{\varphi(n)}$$
        converges in $V$ and its limit does not depend on the choice of $\varphi$.
        We denote the limit by $$\sum_{i\in I} a_i x_i.$$
        \item One has 
        $$ \left\| \sum_{i\in I} a_i x_i\right\|^2 = \sum_{i\in I} \left|a_i\right|^2 \cdot \left\|x_i\right\|^2.$$
        Moreover, for any bijection $\varphi : \NN \rightarrow I'$, and any $y \in V$, the series
        $$ \sum_{n\in \NN} \langle a_{\varphi(n)} x_{\varphi(n)}, y\rangle$$
        converges absolutely  to $\dis \langle \sum_{i\in I} a_i x_i, y\rangle$.
    \end{enumerate}
\end{propositionenv}

\begin{proofenv}
    \quad 
    \begin{enumerate}
        \item For any $\varepsilon \in \QQ$, the set $\{i\in I \mid \left|a_i\right|\left\|x_i\right\| > \varepsilon\}$ is finite, since otherwise, $\dis \sum_{i\in I} \left|a_i\right|^2 \cdot \left\|x_i\right\|^2$ would be infinite.
        Therefore, 
        $$ I' = \bigcup_{\varepsilon \in \QQ} \left\{i\in I \mid \left|a_i\right|\left\|x_i\right\| > \varepsilon\right\}$$
        is countable.
        \item We will prove that the series is a Cauchy sequence, i.e. for any $\varphi : \NN \rightarrow I'$, $\dis \sum_{n\in\NN} a_{\varphi(n)} x_{\varphi(n)}$ is a Cauchy sequence.
        For any $N\in \NN$, $n,m\in \NN$, $n>m\ge N$, one has 
        $$ \left\| \sum_{k = m + 1}^{n} a_{\varphi(k)} x_{\varphi(k)}\right\|^2 = \sum_{k = m + 1}^{n} |a_{\varphi(k)}|^2 \cdot \|x_{\varphi(k)}\|^2.$$
        We assume that
        $$ \sum_{n\in \NN} |a_{\varphi(n)}|^2 \|x_{\varphi(n)}\|^2 = \sum_{i \in I} |a_i|^2 \|x_i\|^2 < +\infty.$$
        We deduce that 
        $$ \lim_{N\rightarrow +\infty} \sup_{n>m\ge N} \left\| \sum_{k = m + 1}^{n} a_{\varphi(k)} x_{\varphi(k)}\right\|^2 = 0.$$
        Hence $\dis  \sum_{n\in \NN} a_{\varphi(n)} x_{\varphi(n)}$ is a Cauchy sequence and converges to $x \in V$.
        
        \quad For any $n\in \NN$, 
        $$ \begin{aligned}
            \left\| x - \sum_{n = 0}^{N} a_{\varphi(n)} x_{\varphi(n)}\right\|^2 &= \lim_{m\rightarrow +\infty} \left\| \sum_{n = 0}^{N + m} a_{\varphi(n)} x_{\varphi(n)} - \sum_{n = 0}^{N} a_{\varphi(n)} x_{\varphi(n)}\right\|^2\\
            & = \lim_{m\rightarrow +\infty} \left\| \sum_{n = N + 1}^{N + m} a_{\varphi(n)} x_{\varphi(n)}\right\|^2\\
            & = \sum_{n = N + 1}^{\infty} |a_{\varphi(n)}|^2 \|x_{\varphi(n)}\|^2
        \end{aligned}$$
        We need to check that $\dis \sum_{ n\in \NN} a_{\varphi(n)} x_{\varphi(n)}$ does not depend on the choice of $\varphi$.

        \quad Let $\psi: \NN \rightarrow I'$ be another bijection. 
        For any $N\in \NN$, there exists $M\in \NN$ such that
        $$ \{\varphi(0),\cdots \varphi(N)\}\subseteq \{\psi(0), \cdots \psi(M) \},\ M \coloneq \max\{\psi^{-1} \left(\varphi(i)\right)\}.$$
        equivalently,
        $$\{\psi(n)\mid  b\ge M + 1 \}\subseteq \{\varphi(m) \mid n\ge N + 1\}.$$
        Therefore, 
        $$ \left\| x - \sum_{n = 0}^{M} a_{\psi(n)} x_{\psi(n)}\right\|^2 = \sum_{n = M + 1}^{\infty} |a_{\psi(n)}|^2 \|x_{\psi(n)}\|^2\le \sum_{n = N + 1}^{+\infty} |a_{\varphi(n)}|^2 \|x_{\varphi(n)}\|^2,$$
        which gives
        $$ \lim_{M\rightarrow +\infty} \left\| x - \sum_{n = 0}^{M} a_{\psi(n)} x_{\psi(n)}\right\|^2 = 0.$$
        \item Pick a bijection $\varphi : \NN\rightarrow I'$, the mapping $\|\cdot\|^2$ is continuous,
        $$ \begin{aligned}
            \left\| \sum_{i\in I} a_i x_i\right\|^2 &= \lim_{N\rightarrow +\infty} \left\| \sum_{n = 0}^{N} a_{\varphi(n)} x_{\varphi(n)}\right\|^2\\
            & = \lim_{N\rightarrow +\infty} \sum_{n = 0}^{N} \left|a_{\varphi(n)}\right|^2 \cdot \left\|x_{\varphi(n)}\right\|^2\\
            & = \sum_{i\in I} \left|a_i\right|^2 \cdot \left\|x_i\right\|^2
        \end{aligned}$$
        By Cauchy Schwarz,
        $$\begin{aligned}
            \left| \langle \sum_{i\in I} a_i x_i, y\rangle - \sum_{n=0}^{N}\langle a_{\varphi(n)} x_{\varphi(n)}, y\rangle \right| = & \left| \left\langle \sum_{n =N+1}^{+\infty} a_{\varphi(n)} x_{\varphi(n)}, y\right\rangle \right|\\
            & \le \left(\sum_{n=N+1}^{+\infty} |a_{\varphi(n)}|^2 \right)^{1/2} \cdot \left\|y\right\|.
        \end{aligned}$$
        Hence, 
        $$\lim_{N\rightarrow +\infty} \left| \langle \sum_{i\in I} a_i x_i, y\rangle - \sum_{n=0}^{N}\langle a_{\varphi(n)} x_{\varphi(n)}, y\rangle \right| = 0.$$
        Now, we will check that $\dis i\in I \left| \left\langle a_i x_i, y\right\rangle\right|<+\infty $.

        \quad For any $i\in I$, pick $b_i \in K$, such that $|b_i|=1$,
        $$ \langle b_i a_i x_i, y\rangle = \left|\langle a_i x_i, y\rangle\right|.$$
        We know that 
        $$ \sum_{i\in I} |b_i a_i|^2 \left\|x_i\right\|^2 = \sum_{i\in I} \left|a_i\right|^2 \cdot \left\|x_i\right\|^2 < +\infty.$$
        We obtain,
        $$ \sum_{i\in I} |\langle a_i x_i, y\rangle|^2 = \sum_{n\in \NN} \langle b_{\varphi(n)} a_{\varphi(n)} x_{\varphi(n)}, y\rangle^2 < +\infty.$$
    \end{enumerate}
\end{proofenv}
\begin{propositionenv}[Parseval]
    Let $\left(V,\langle \cdot, \cdot\rangle\right)$ be an inner product space, $W\subseteq V$ be a subspace, such that $\left(W,\langle\cdot,\cdot\rangle|_W\right)$ is a Hilbert space.
    \begin{enumerate}
        \item The mapping $p_W : V\rightarrow W$ is $K$-linear.
        \item For any $x\in W$, $p_W(x) = x$. In particular, $p_W\circ p_W = p_W$.
        \item Let $\left(e_i\right)_{i\in I}$ be an orthonormal basis of $W$, then for any $x\in V$,
        $$ \sum_{i\in I} |\langle e_i, x\rangle|^2 \le \| x\|^2,$$
        and 
        $$ p_W(x) = \sum_{i\in I} \langle e_i, x\rangle e_i.$$
    \end{enumerate}
\end{propositionenv}
\begin{proofenv}
    \quad
    \begin{enumerate}
        \item Let $(x,y)\in V\times V$, $(a,b)\in K^2$. For any $z\in W$,
        $$ \langle z,\left(ax +by\right) - \left(a p_W(x) + bp_W(y)\right)\rangle = a\langle z, x - p_W(x)\rangle + b\langle z, y - p_W(y)\rangle,$$
        which is $0$. By theorem \ref{3.2.2} we obtain,
        $$p_W(x) (ax + by) = ap_W(x) + bp_W(x).$$
        \item For any $z\in W$, $0 = z-z$ belongs to $W^\perp$. Hence, $p_W(z) = z$.
        We also get $p_W(p_W(x)) = p_W(x)$, $\forall x\in V$.
        \item First consider finite $I=\{1,\cdots,n\}$. Take
        $$ y\coloneq x - \sum_{i=1}^{n} \langle e_i , x\rangle e_i.$$
        For any $j \in \{1,\cdots ,n\}$,
        $$ \langle e_j,y \rangle = \langle e_j, x\rangle - \langle e_j, x\rangle \langle e_j,e_j\rangle = 0.$$
        So, $y\in W^\perp$. Therefore $\dis p_W(x) = \sum_{i=1}^{n} \langle e_i,x\rangle e_i.$
        In particular,
        $$ \left\| p_W(x)\right\|^2 = \sum_{i=1}^{n} \left|\langle e_i,x\rangle\right|^2 \le \left\|x\right\|^2.$$
        Now $I$ is general. For any $J\subseteq I$ finite,
        $$ \sum_{i\in J} |\langle e_i,x\rangle|^2 \le \|x\|^2.$$
        Taking sup over all possible $J$, we obtain
        $$ \sum_{i\in I} |\langle e_i,x\rangle|^2 \le \|x\|^2.$$
        So $\dis \sum_{i\in I} |\langle e_i,x\rangle|^2 $ is well-defined.
        For an $j\in I$, 
        $$ \langle e_j,\sum_{i\in I} |\langle e_i,x\rangle e_i\rangle = \langle e_i,x\rangle.$$
        So, 
        $$ x - \sum_{i\in I} \langle e_i, x\rangle e_i \in \{e_i \mid i\in I\}^\perp = W^\perp$$
        which allows us to conclude that $\dis p_W(x) = \sum_{i\in I} \langle e_i,x\rangle.$
    \end{enumerate}
\end{proofenv}

\begin{theoremenv}[Gram-Schmidt]
    Let $\left(V,\langle \cdot, \cdot\rangle\right)$ be an inner product space,
    $$ \{0\} = V_0\subseteq V_1 \subseteq \cdots \subseteq V_n$$
    sequence of vector space such that $\forall i\in \{1,\cdots,n\}$, $\dim \left(V_i\right) = i.$
    Then there exists an orthonormal family $\left(x_i\right)_{i=1}^n$ in $V$ such that $x_{i} \in V_i\setminus V_{i-1}$ for any $i\in \{1,\cdots,n\}$.
\end{theoremenv}
\begin{proofenv}
    For any $i \in \{1,\cdots,n\}$, pick $y_{i}\in V_i\setminus V_{i-1}$ and 
    $$ x_i = \frac{y_i - p_{V_i} (y_i)}{\| y_i - p_{V_{i-1}}(y_i)\|},\ x_i \in \left(V_i \setminus V_{i-1}\right)\cap V_{i-1}^\perp.$$
    $\left(x_i\right)_{i=1}^n$ is orthonormal.
\end{proofenv}
\begin{remark}
    In the proof of Gram-Schmidt theorem, the family $\left(x_j\right)_{j=1}^{i-1}$ forms an orthogonal basis.
    Hence the vector $x_i$ can be construct in recursive way:
    $$ z_i \coloneq y_i - \sum_{j=1}^{i-1} \langle y_i,x_j\rangle x_j, y_i \in V_i\setminus V_{i-1},\ x_i = \frac{z_i}{\|z_i\|}.$$
    $$ \|z_i\|^2 = \|y_i\|^2 - \sum_{j=1}^{i-1} |\langle x_j ,y_j \rangle|^2 = \|y_i\|^2 - \sum_{j=1}^{i-1}\frac{|\langle z_i, y_i\rangle|^2}{\|z_j\|^2}.$$
\end{remark}


\section{Adjoint of a Bounded Linear Mapping}

\begin{lemmaenv}\label{lem:3.3.1}
    Let $\left(V,\langle \cdot,\cdot\rangle\right)$ be a Hilbert space and $W\subseteq V$ be a vector space of $V$.
    If $W$ is not dense in $V$, then $W^\perp$ is non-zero.
\end{lemmaenv}
\begin{proofenv}
    By proposition \ref{prop:3.1.9}, if $W' = \overline{W}$, then $\left(W'\right)^\perp = W^\perp$, so without loss of generality,
    we can assume that $W$ is closed (i.e. $\left(W, \left.\langle\cdot,\cdot\rangle\right|_{W}\right)$ is a Hilbert space) and $W \neq V$.
    Let $x\in V\setminus W$. Then $p_W (x)\neq x$.
    $$ x - p_W(x) \in W^\perp\setminus \{0\}.$$
\end{proofenv}

\begin{definitionenv}
    Let $V,W$ be vector spaces over $K$.
    We call \textbf{conjugate homomorphism} from $V$ to $W$ any mapping $f:V\rightarrow W$ such that 
    for any $(a,b)\in K^2$ and $(x,y)\in V^2$,
    $$ f(ax +by) = \bar{a} f(x) + \bar{b}f(y).$$
    If in addition, $f$ if a bijection, we say that $f$ is a \textbf{conjugate isomorphism}.
\end{definitionenv}

\begin{definitionenv}
    Let $\left(V, \langle \cdot,\cdot\rangle\right)$ be a Hilbert space.
    We say that a $K$-linear mapping $f : V\rightarrow V$ is \textbf{isometric} (or \textbf{isometry}) if 
    $$ \| f(x)\| = \|x\|, \ \forall x\in V.$$
    Note that isometry is automatically bounded and injective.
\end{definitionenv}

\begin{theoremenv}[Riesz Representation Theorem]
    Let $\left(V, \langle \cdot,\cdot\rangle\right)$ be a Hilbert space.
    For any $f: V\rightarrow K$ that is linear and bounded, there exists a \textit{unique} $x\in V$ such that
    $$ f = l_x \coloneq \langle x,\cdot\rangle.$$
    The mapping 
    $$\begin{array}{rcl}
        V & \rightarrow & \mathcal{L}(V,K)\\
        x & \mapsto & l_x
    \end{array}$$
    is a conjugate isomorphism, that is also an isometry, when we consider $\mathcal{L}(V,K)$ with the standard operator norm.
\end{theoremenv}

\begin{proofenv}
    By proposition \ref{prop:3.1.8}, if the mapping $V\rightarrow \mathcal{L}(V,K)$, $x\mapsto l_x$ is well-defined, injective and for any $x\in V$, $\| l_x\| = \| x\|_{V}$.
    Let $(x,y)\in V^2$, $(a,b)\in K^2$,
    $$ \forall z \in V, l_{ax + by} (z) = \langle ax+by, z\rangle = \bar{a} \langle x,z\rangle + \bar{b}\langle y,z\rangle = \bar{a} l_x(z) + \bar{b} l_y(z).$$
    So the mapping is a conjugate homomorphism.
    It remains to show that the mapping $x\mapsto l_x$ is surjective.

    \quad $f$ is bounded and linear (hence continuous) so $\ker f\eqcolon W$ is a closed vector subspace of $V$. 
    If $f$ is the zero mapping, then $x = 0$ is the only element of $V$ such that $l_x = f$.
    For non-zero $f$, $\overline{W}\neq V$, hence by lemma \ref{lem:3.3.1}, there exists $0\neq y\in W^\perp$ such that $\|y\| =1$.
    Put $x \coloneq \overline{f(y)}y$.
    We have $V/W \cong K$.
    So for any $z\in V$, there exists $\lambda \in K$ such that $z - \lambda y\in W$.
    $$ l_x(z) = \langle \overline{f(y)} y,z\rangle = f(y) \langle y,z\rangle = f(y) \langle y,\lambda y\rangle = \lambda f(y) = f(\lambda y) = f(z).$$
\end{proofenv}

\begin{corollaryenv}
    Let $\left(V, \langle \cdot,\cdot\rangle_V\right)$ and $\left(W, \langle \cdot,\cdot\rangle_W\right)$ be Hilbert spaces
    and $f: V\rightarrow W$ be a bounded linear mapping.
    There exists a \textit{unique} bounded linear mapping $f^*: W \rightarrow V$ such that 
    $$ \forall (x,y)\in W\times V, \langle f^*(x), y\rangle_V = \langle x, f(y) \rangle_W.$$
    The mapping $f^*$ is called \textbf{adjoint} of $f$.
    The operator norm of $f^*$ is equal to that of $f$, and one has
    $$ \left(f^*\right)^* = f.$$
\end{corollaryenv}
\begin{proofenv}
    For any $x\in W$, the mapping $l_x \circ f : V\rightarrow K$ is linear bounded.
    By Riesz representation theorem (applied to $l_x \circ f$), there exists a unique element of $f$, which we denote by $f^*(x)$, such that
    $$ l_x\circ f = l_{f^*(x)}.$$
    For any $(x,x')\in W^2$ and $(a,b)\in K^2$, 
    $$ \begin{aligned}
        l_{f^*(ax+bx')} &= l_{ax+bx'}\circ f = \left(\bar{a}l_x + \bar{b} l_{x'}\right)\circ f = \bar{a}\left(l_x \circ f\right) + \bar{b} \left(l_{x'} \circ f\right)\\
        &= \bar{a} l_{f^*(x)} + \bar{b} l_{f^*(x')} = l_{a f^*(x)+bf^*(x')}.
    \end{aligned}$$
    So $f^*: W\rightarrow V$ is unique linear mapping such that 
    $$ \forall (x,y)\in W\times V,\ \langle f^*(x), y\rangle_V = \langle x, f(y) \rangle_W.$$
    Let's check that $f^*$ is bounded.
    $$\forall x\in W,\ \| f^*(x)\|_V =\| l_{f^*(x)}\|  = \| l_x\circ f\| \le \|l_x\| \|f\| = \|x\|_{W} \|f\|.$$
    Hence $\| f^*\|\le \|f\|$. For any $(x,y)\in W\times V$,
    $$ \langle y, f^*(x)\rangle_V = \overline{\langle f^*(x), y\rangle_{V}} = \overline{\langle x, f(y)\rangle_W} = \langle f(x), x\rangle_W.$$
    So $\left(f^*\right)^* = f$, which also leads to 
    $$ \| f\| = \| \left(f^*\right) ^* \| \le \|f^*\|.$$
    So, $\|f\| = \|f^*\|$.
\end{proofenv}

\begin{exampleenv}
    Let $n$, $p$ be two positive integers, and $\langle \cdot,\cdot\rangle_{n}$, $\langle \cdot, \cdot\rangle_{p}$ be the canonical inner product on $K^n$ and $K^p$ respectively.
    Let $A:K^n \rightarrow K^p$ be $K$-linear. $A$ can be represented as
    $$ A \bumpeq  \begin{pmatrix}
        a_{1,1} & \cdots & a_{1,p}\\
        \vdots & \ddots & \vdots\\
        a_{n,1} & \cdots & a_{n,p}
    \end{pmatrix}.$$
    $$ A: (y_1, \cdots, y_n)\longmapsto \left(\sum_{j=1}^{n} y_j a_{j,1}, \cdots,\sum_{j=1}^{n}y_j a_{j,p}\right).$$
    For any $x= (x_1,\cdots,x_p)\in K^p$, the mapping $l_x\circ A$ works the following way:
    $$ (y_1,\cdots y_n) \mapsto \left(\sum_{i=1}^{p}\bar{x_i} \sum_{j=1}^{n}y_ja_{j,i}\right)$$
    Right hand side is $\dis \sum_{j=1}^{n}\overline{\left(\sum_{i=1}^{p} x_i \overline{a_{j,i}}\right)}y$. So,
    $$ A^* \bumpeq \begin{pmatrix}
        \overline{a_{1,1}} & \cdots & \overline{a_{1,p}}\\
        \vdots & \ddots & \vdots\\
        \overline{a_{n,1}} & \cdots & \overline{a_{n,p}}
    \end{pmatrix}.$$
\end{exampleenv}

\begin{propositionenv}
    Let $\left(V, \langle \cdot,\cdot\rangle_{V}\right)$ and $\left(W, \langle \cdot,\cdot\rangle_{W}\right)$ be Hilbert spaces.
    Then $f\mapsto f^*$ defines a conjugate homomorphism from $\mathcal{L}(V,W)$ to $\mathcal{L}(W,V)$.
\end{propositionenv}
\begin{proofenv}
    Let $(f,g) \in \mathcal{L}(V,W)\times \mathcal{L}(V,W)$, $(a,b) \in K\times K$.
    For any $(x,y)\in W\times V$, 
    $$\begin{aligned}
        \langle x,(af+bg)(y) \rangle_{W} & =  a\langle x,f(y)\rangle_{W} + b\langle x,g(y)\rangle_{W}\\
        &= a\langle f^*(x),y\rangle_V + b\langle g^*(x),y\rangle_V \\
        &= \langle \bar{a} f^*(x) + \bar{b} g^*(x), y\rangle_V\\
        &= \langle (af+bg)^*(x), y\rangle_V.
    \end{aligned}$$
    So $(af + bg)^* = \bar{a} f^* + \bar{b} g^*$.
\end{proofenv}

\begin{propositionenv}
    Let $(V_1,\langle\cdot,\cdot\rangle_{1})$, $(V_2,\langle\cdot,\cdot\rangle_{2})$, $(V_3,\langle\cdot,\cdot\rangle_{3})$ be Hilbert spaces,
    $f: V_1\rightarrow V_2$, $g:V_2\rightarrow V_3$ be bounded linear mappings. 
    Then 
    $$ (g\circ f)^* = f^*\circ g^*.$$
\end{propositionenv}
\begin{proofenv}
    For $(x,y) \in V_3\times V_1$,
    $$\begin{aligned}
        \langle x, (g\circ f)(y)\rangle_{3} &= \langle x,g(f(y))\rangle_{3} = \langle g^*(x), f(y)\rangle_{2}\\
        &= \langle f^*(g^*(x)), y\rangle_{1} = \langle (f^*\circ g^*)(x), y\rangle_{1}.
    \end{aligned}$$
    So, $\left(g\circ f\right)^* = f^*\circ g^*$.
\end{proofenv}

\begin{definitionenv}
    Let $\left(V,\langle\cdot,\cdot\rangle\right)$ be a Hilbert space.
    We say that a bounded endomorphism $f:V\rightarrow V$ is self-adjoint if $f^* = f$:
    $$ \forall (x,y)\in V\times V, \ \langle x,f(y)\rangle = \langle f(x),y\rangle.$$
\end{definitionenv}

\begin{definitionenv}
    Let $n\in \NN_{>0}$.
    We equipped $K^n$ with the canonical inner product
    $$ \langle x,y\rangle = \sum_{i=1}^{n} \overline{x}_i y_i.$$
    Let $A:K^n\rightarrow K^n$ be linear, then $A^* = \overline{A^{\mathsf{T}}}$.

    \quad In the cases when $K =\RR$ and $A$ is self-adjoint (i.e. $A = A^* = A^{\mathsf{T}}$), we say that the matrix of $A$ is a \textbf{symmetric matrix}.

    \quad In the case when $K = \CC$ and $A$ is self-adjoint, (i.e. $A = A^* = \overline{A^{\mathsf{T}}}$), we say that the matrix of $A$ is a \textbf{Hermitian matrix}.
\end{definitionenv}

\begin{exampleenv}
    Let $(V,\langle \cdot,\cdot\rangle)$ be a Hilbert space and $W$ be a closed vector subspace of $V$.
    Then the mapping of orthogonal projection $p_W:V\rightarrow V$ is self-adjoint.
    
    \quad Indeed, for any $(x,y)\in V\times V$, one has
    $$ \langle x, p_W(y)\rangle = \langle p_W(x), p_W(y)\rangle + \langle x - p_W(x), p_W(y)\rangle = \langle p_W(x), p_W(y)\rangle,$$
    since $x - p_W(x) \in W^\perp$ and $p_W(y)\in W$.
    Similarly, 
    $$ \langle p_W(x), y\rangle = \langle p_W(x), p_W(y)\rangle + \langle p_W(x), y - p_W(y)\rangle = \langle p_W(x), p_W(y)\rangle.$$
\end{exampleenv}



\section{Reduction of Self-adjoint Mappings}

\begin{definitionenv}
    Let $F$ be a field, $V$ be a vector space over $F$ and $f:V\rightarrow V$ be an endomorphism.
    Let $\lambda\in F$.
    If there exists $x\in V\setminus\{0\}$ such that
    $$f(x) = \lambda x,$$
    then we say that $\lambda$ is an \textbf{eigenvalue} of $f$ and that $x$ is an \textbf{eigenvector} of $f$ associated with the eigenvalue $\lambda$.
    
    \quad Note that $\lambda \in F$ is an eigenvalue of $f$ if and only if $\lambda \Id_{V} - f$ is not injective.
    Suppose that $V$ is finitely-dimensional over $F$. 
    We say that the endomorphism $f: V\rightarrow V$ is \textbf{diagonalizable} if there exists a basis of $V$ composed of eigenvectors of $f$.
\end{definitionenv}

\begin{remark}
    Let $F$ be a field, $V$ be a finitely-dimensional vector space over $F$, $f:V\rightarrow V$ be an endomorphism.
    Let $\left(e_i\right)_{i=1}^{n}$ ($n = \dim_F(V)$) be a basis such that $\forall i\in \{1,\cdots,n\}$, $e_i$ is an eigenvector of $f$.
    For any $i\in \{1,\cdots,n\}$, let $\lambda_i \in F$ be such that 
    $$ f(e_i) = \lambda_i e_i.$$
    Let
    $$\begin{array}{rrcl}
        \varphi: & K^n & \rightarrow & V\\
        & (a_1,\cdots,a_n) & \mapsto & \sum_{i=1}^{n} a_i e_i
    \end{array}$$
    be a natural isomorphism $K^n \cong V$, then 
    $$ \begin{array}{rrcl}
        \varphi^{-1} \circ f \circ \varphi: & K^n & \rightarrow & K^n\\
        & (a_1,\cdots,a_n) & \mapsto & (\lambda_1 a_1,\cdots,\lambda_n a_n)
    \end{array}$$
    i.e. $\varphi^{-1}\circ f \circ \varphi = \mathrm{diag} (\lambda_1,\cdots,\lambda_n)$.
\end{remark}

\begin{definitionenv}
    Let $(V,\|\cdot\|)$ be a Banach space over $K$ ($K = \RR $ or $\CC$), $f:V\rightarrow V$ bounded linear mapping.
    Note that the \textit{open mapping theorem} shows that if $\lambda \Id_{V}-f$ is a bijection, then $\left(\lambda\Id_{V} - f\right)^{-1}$ is bounded, linear.
    We denote by $\sigma (f)$ the set of $\lambda \in K$ such that $\lambda\Id_{V} - f$ is invertible.
    We call $\sigma(f)$ the \textbf{resolvent of $f$}.

    The set 
    $$ \rho(f) = K\setminus \sigma(f)$$
    is called the \textbf{spectrum} of $f$.
    
    \quad In the case when $V$ is finite dimensional, a $K$-linear mapping is an isomorphism if and only if it is injective.
    Therefore, in this case, the spectrum of $f$ is equal to the set of eigenvalue of $f$.
    {\bf This is not true in general!}

    \quad In the infinite dimensional case, the points of the spectrum can be divided into two classes:
    \begin{enumerate}[label= \Roman*.]
        \item $\lambda \in \rho(f)$ such that $\lambda \Id_{V} - f$ is not injective.
        \item $\lambda \in \rho(f)$ such that $\lambda \Id_{V} - f$ is injective but not surjective.
    \end{enumerate}
    The points of type I correspond to the eigenvalues of $f$ and the set of points of type II might be non-empty.
\end{definitionenv}

\begin{exampleenv}
    Consider a Hilbert space of countable dimensional over $K$ with an orthonormal basis $\{e_n\}_{n\in\NN}$ and let $f$ be such that $f(e_n) = \frac{1}{n} e_n$ for all $n\in \NN$.
    The set of eigenvalues of $f$ if exactly $\{\frac{1}{n} \mid n\in \NN\}$ while the spectrum contains also the value $0\in K$.
    Indeed, $f^{-1}$ sends $e_n$ to $ne_n$ for all $n\in \NN$ and therefore is not bounded.
    The open mapping theorem implies that $f$ cannot be surjective.
    This implies that $0\in \rho(f)$ and $0$ is not an eigenvalue of $f$.
    ($0\in K$ is a point of type II)
\end{exampleenv}
\begin{theoremenv}
    Let $\left(V, \langle\cdot,\cdot\rangle\right)$ be a finitely dimensional inner product space and let $f:V\rightarrow V$ be a self-adjoint endomorphism.
    There exists an orthonormal basis of $V$ composed of eigenvectors of $f$.
    
    \quad In particular, $f$ is diagonalizable.
    Moreover, the eigenvalues of $f$ must be all real. 
\end{theoremenv}
\begin{proofenv}
    If $\lambda$ is an eigenvalue of $f$ and $v\in V$ is an eigenvector of the eigenvalue, then
    $$ \lambda \langle v,v\rangle = \langle v,\lambda v\rangle = \langle v,f(v)\rangle = \langle f(v), v\rangle = \bar{\lambda} \langle v,v\rangle.$$
    So $\lambda \in \RR$.

    \quad We reason by induction on $\dim_{K} (V) = n$.
    When $n = 0$, the statement is trivial.
    Suppose $n = \dim_{K} V >0$, since $f$ is self-adjoint
    $$ \langle x,f(x)\rangle = \langle f(x),x \rangle = \overline{\langle x, f(x)\rangle}$$
    for all $x \in V$.
    Hence $\langle x,f(x)\rangle \in \RR$ for all $x \in V$.
    Define 
    $$ \begin{array}{rrcl}
        q: & V & \rightarrow & \RR\\
        & x & \mapsto & \langle x,f(x)\rangle.
    \end{array}
    $$
    Let $B\coloneq \{x \in V \mid \|x\| \le 1\}$, $S\coloneq \{x \in V\mid \|x\| = 1\}$.
    $S$ is closed bounded in finitely-dimensional vector space $V$, so it is compact.
    The mapping $q|_S$ attains its maximum $\lambda$.
    For any $y \in V$, there exists $a\in K$ and $x \in S$ such that $y = a x$.
    $$ q(y) = \langle ax, f(ax)\rangle = |a|^2 \langle x,f(x)\rangle = |a|^2 q(x) \le |a|^2 \lambda = \lambda\|y\|^2.$$
    Let $e_1$ be an element of $S$ such that $q(e_1) = \lambda$.
    Let $y\in S$ and $\theta \in K$ be such that $|\theta|=1$.
    For any $t \in \RR_{>0}$,
    $$\begin{aligned}
        q(\theta e_1 + ty) = \langle \theta e_1 + ty, f(\theta e_1 + ty)\rangle &= q(e_1) + t^2 q(y) + 2t\mathrm{Re}\left(\theta\langle y,f(e_1)\rangle\right)\\
        &\le \lambda\|\theta e_1 + ty\|^2,
    \end{aligned}$$
    $$ \lambda \|\theta e_1 + ty\|^2 = \lambda\left(1+ t^2 +2t \mathrm{Re}(\theta\langle y,e_1\rangle)\right),$$
    so
    $$ tq(y) + 2\mathrm{Re} \left(\theta \langle y,f(e_1)\rangle\right)\le \lambda t + 2\lambda \mathrm{Re} \left(\theta\langle y,e_1\rangle\right).$$
    Taking the limit as $t\rightarrow 0$, we obtain
    $$ \mathrm{Re} \left(\theta \langle y,f(e_1)\rangle\right) \le \lambda \mathrm{Re} \left(\theta \langle y,e_1\rangle\right).$$
    Replacing $y$ with $-y$, we get
    $$ \mathrm{Re} \left(\theta \langle y,f(e_1)\rangle\right)\ge \lambda \mathrm{Re} \left(\theta \langle y,e_1\rangle\right).$$
    We deduce that $\mathrm{Re}\left(\theta\langle y,f(e_1)\rangle\right) = \lambda \mathrm{Re}\left(\theta \langle y, e_1\rangle\right)$, so $\mathrm{Re}\left(\theta\langle y, f(e_1)-\lambda e_1\rangle\right) = 0.$
    Since $\theta$ is arbitrary but $|\theta|=1$, we can fix it in a way that
    $$\mathrm{Re}\left(\theta\langle y, f(e_1) - \lambda e_1\rangle\right) = | \langle y, f(e_1) - \lambda e_1\rangle|.$$
    Thus $f(e_1) = \lambda e_1$, i.e. $\lambda$ is an eigenvalue of $f$ and $e_1$ is its eigenvector.
    Let $W = \{e_1\}^{\perp}$, $\dim_{K}(W) = \ker(l_{e_1}) = n - 1$.
    $$ \langle e_1, f(y) \rangle = \langle f(e_1) ,y\rangle = \langle \lambda e_1 ,y \rangle = \bar{\lambda} \langle e_1, y\rangle = 0.$$
    The restriction of $f$ to $W$ gives a self-adjoint operator $f|_{W} : W\rightarrow W$.
    After using the inductive assumption, we obtain an orthonormal basis $\left(e_i\right)_{i=1}^{n}$ of $W$ composed of eigenvectors of $f|_{W}$.
    Since $e_1\in W^\perp$ and $\|e_1\| =1$, we conclude that $\left(e_i\right)_{i=1}^{n}$ is the desired basis of $V$.
\end{proofenv}

\begin{corollaryenv}
    Let $(V,\langle\cdot,\cdot\rangle)$ be a finite-dimensional inner product space and $f:V\rightarrow V$ be a self-adjoint $K$-linear mapping (bounded since $\dim_{K}(V) < \infty$).
    The operator norm of $f$ is equal to the maximal eigenvalue of $f$.
\end{corollaryenv}
\begin{proofenv}
    Let $\lambda \in K$ be an eigenvalue and $v$ be its eigenvector that $\|v\| = 1$.
    By Cauchy-Schwarz,
    $$ |\lambda| = |\langle v,f(v)\rangle| \le \| v\| \cdot\|f(v)\| \le \|f\|.$$
    Let $(e_i)_{i=1}^{n}$ be an orthonormal basis of $V$ composed of eigenvectors.
    Let $x = a_1 e_1 + \cdots + a_n e_n$, with $a_i \in K$ be an arbitrary element of $V$,
    $$ f(x) = a_1 f(e_1) + \cdots + a_n f(e_n) = a_1 \lambda_1 e_1 + \cdots + a_n \lambda_n e_n.$$
    So, 
    $$\begin{aligned}
        \|f(x)\|^2 = |a_1 \lambda_1|^2 + \cdots + |a_n \lambda_n|^2 &\le \max_{i\in \{1,\cdots,n\}} |\lambda_i|^2 \left(|a_1|^2 + \cdots + |a_n|^2\right)\\
        &= \max_{i\in \{1,\cdots,n\}} |\lambda_i|^2 \|x\|^2.
    \end{aligned}$$
    Thus, $\|f\|^2 \le \max_{i\in \{1,\cdots,n\}} |\lambda_i|^2.$
\end{proofenv}

\begin{definitionenv}
    Let $\left(V,\langle\cdot,\cdot\rangle\right)$ be a Hilbert space, and $f:V\rightarrow V$ be a self-adjoint bounded $K$-linear mapping.
    We say that $f$ is 
    \begin{enumerate}[label=\roman*]
        \item positive semi-definite if $\langle f(x),x\rangle \ge 0$ for all $x\in V$.
        \item negative semi-definite if $\langle f(x),x\rangle \le 0$ for all $x\in V$.
        \item positive definite if $\langle f(x),x\rangle > 0$ for all $x\in V\setminus \{0\}$.
        \item negative definite if $\langle f(x),x\rangle < 0$ for all $x\in V\setminus \{0\}$.
    \end{enumerate}
\end{definitionenv}

\begin{propositionenv}
    Let $(V,\langle\cdot,\cdot\rangle)$ be a finite-dimensional inner product space and $f:V\rightarrow V$ be a self-adjoint $K$-linear mapping.
    \begin{enumerate}
        \item $f$ is positive semi-definite if and only if all eigenvalues of $f$ are non-negative.
        \item $f$ is negative semi-definite if and only if all eigenvalues of $f$ are non-positive.
        \item $f$ is positive definite if and only if all eigenvalues of $f$ are strictly positive.
        \item $f$ is negative definite if and only if all eigenvalues of $f$ are strictly negative.
    \end{enumerate}
\end{propositionenv}
\begin{proofenv}
    Let $(e_i)_{i=1}^{n}$ be an orthonormal basis of $V$ composed of eigenvectors of $f$.
    For $i\in \{1,\cdots,n\}$, let $\lambda_i$ be the eigenvalue associated with $e_i$.
    If $f$ is positive semi-definite (resp. semi-definite, etc.), then for any $i\in \{1,\cdots,n\}$,
    $$ \lambda_i = \langle f(e_i), e_i\rangle \ge 0.$$
    Conversely, for any $x = a_1e_1 + \cdots + a_n e_n \in V\in \{0\}$ one has
    $$\begin{aligned}
        \langle f(x),x\rangle &= \langle a_1 e_1 + \cdots + a_n e_n, f(a_1 e_1 + \cdots + a_n e_n)\rangle\\
        &= \langle a_1 e_1 + \cdots + a_n e_n, a_1 \lambda_1 e_1 + \cdots + a_n \lambda_n e_n\rangle\\
        &= \lambda_1 |a_1|^2 + \cdots + \lambda_n |a_n|^2.
    \end{aligned}$$
    If for any $i \in \{1,\cdots,n\}$, $\lambda_i\ge 0$, then one has that $\langle x, f(x) \rangle\ge 0$ as required.
\end{proofenv}


\section{Linear Isometry}

\begin{propositionenv}
    Let $(V,\langle\cdot,\cdot\rangle)$ be a Hilbert space.
    A $K$-linear mapping $f : V\rightarrow V$ is isometric if and only if
    $$ f^*\circ f = \Id_V,$$
    or equivalently, 
    $$ \forall x,y\in V, \ \langle f(x),f(y)\rangle = \langle x,y\rangle.$$
\end{propositionenv}
\begin{proofenv}
    For any $(x,y)\in V\times V$,
    $$ \langle f(x),f(y)\rangle = \langle f^*(f(x)),y\rangle .$$
    Suppose that $f^*\circ f = \Id_V$, then for any $(x,y)\in V\times V$,
    $$ \langle f(x),f(y)\rangle = \langle f^*(f(x)),y\rangle = \langle x,y\rangle.$$
    In particular, for any $x\in V$, 
    $$ \| f(x)\|^2 = \langle f(x),f(x)\rangle = \langle x,x\rangle,$$
    so $f$ is isometric.

    \quad Conversely, suppose that $f$ is isometric.
    Let $(x,y)\in V\times V$, $\theta \in K$ such that $|\theta| = 1$.
    One has
    $$ \|f(x+\theta y)\|^2 = \|f(x)\|^2 + \|f(y)\|^2 + 2 \mathrm{Re} \left(\theta \langle f(x),f(y)\rangle\right),$$
    $$ \| x+ \theta y\|^2 = \|x\|^2 + \|y \|^2 + 2 \mathrm{Re} \left(\theta \langle x,y\rangle\right),$$
    so 
    $$ \mathrm{Re} \left(\theta \langle f^*(f(x)),y \rangle\right) = \mathrm{Re}(\theta\langle x,y\rangle),\ \mathrm{Re} \left(\theta \langle f^*(f(x)) - x,y\rangle\right) = 0.$$
    Fix $\theta\in K$ in a way that $\theta \langle f^*(f(x)) - x,y\rangle \in \RR$.
    We deduce that $f^*(f(x)) = x$.
\end{proofenv}

\begin{remark}
    Let $(V,\langle\cdot,\cdot\rangle)$ be a Hilbert space,
    $f: V\rightarrow V$ be a isometric $K$-linear mapping.
    Then $f$ must be injective (otherwise $f(x) = f(y)$ for some $x\neq y$ and $\| f(x) - f(y)\| = 0 \neq \|x - y\|$).
    If $\dim V < +\infty$, $f$ is an isomorphism.
    
    \quad In general, an isometric $K$-linear mapping is called an automorphism of the Hilbert space $\left(V, \langle\cdot, \cdot\rangle\right)$.
    The set of automorphisms of $\left(V, \langle\cdot, \cdot\rangle\right)$ forms a group under the composition of mappings, denoted as $\mathrm{Aut}\left(V, \langle\cdot, \cdot\rangle\right)$.
    
    \quad When $K = \RR$, $O(V,\langle\cdot,\cdot\rangle)\coloneq\mathrm{Aut}_{\RR} \left(V, \langle\cdot, \cdot\rangle\right)$ the orthogonal group of $V$.
    
    \quad When $K = \CC$, $U(V,\langle\cdot,\cdot\rangle)\coloneq\mathrm{Aut}_{\CC} \left(V, \langle\cdot, \cdot\rangle\right)$ the unitary group of $V$.
\end{remark}

\begin{exampleenv}
    Let $(V,\langle\cdot,\cdot\rangle)$ be a Hilbert space.
    Let $W \subseteq V$ be a closed $K$-vector subspace.
    We call \textbf{orthogonal symmetry} with respect to $W$ the following mapping:
    $$ \begin{array}{rrcl}
        s_W: & V & \rightarrow & V\\
        & x & \mapsto & 2p_W(x) - x.
    \end{array}$$
    Clearly, $s_W$ is $K$-linear.
    Moreover, for any $x\in V$,
    $$ s_W\left(s_W (x)\right) = 2p_W\left(s_W(x)\right) - s_W(x) = x.$$
    So $s_W$ is an involution ($s_W\circ s_W = \Id_V$).
    Also, 
    $$ s_W^* = 2p_W^* - \Id_V^* = 2p_W - \Id_V.$$
    So $s_W$ is self-adjoint endomorphism of the Hilbert space $\left(V,\langle \cdot,\cdot\rangle\right)$.
\end{exampleenv}
